\textbf{Phần 1. Trắc nghiệm lựa chọn}
\setcounterpageref{ex}{0}

\begin{ex} %[1D6V22]
Trong các số: $45, 54, 49, 51$, số lớn nhất là số nào?
\choice
{$53$}
{$55$}
{$52$}
{\True $54$}
\loigiai{
So sánh chữ số hàng chục: $5 > 4$. Hai số có hàng chục là 5 là 54 và 51, trong đó $54 > 51$. Vậy số lớn nhất là $54$.
\textbf{Đáp số:} $54$.\par
\textbf{Đáp án của câu hỏi là D}
}
\end{ex}

\begin{ex} %[1D6V22]
Trong các số: $78, 87, 69, 90$, số bé nhất là số nào?
\choice
{$68$}
{\True $69$}
{$70$}
{$67$}
\loigiai{
Chữ số hàng chục nhỏ nhất là 6 (trong số 69). Vậy số bé nhất là $69$.
\textbf{Đáp số:} $69$.\par
\textbf{Đáp án của câu hỏi là B}
}
\end{ex}

\begin{ex} %[1D6V22]
Tìm số tự nhiên $x$ thỏa mãn: $38 < x < 40$.
\choice
{$38$}
{$40$}
{\True $39$}
{$37$}
\loigiai{
Số tự nhiên nằm giữa 38 và 40 là số $39$.
\textbf{Đáp số:} $39$.\par
\textbf{Đáp án của câu hỏi là C}
}
\end{ex}

\begin{ex} %[1D6V22]
Cho các số: $23, 67, 45, 89, 12$. Sắp xếp theo thứ tự từ lớn đến bé, số đứng ở vị trí thứ ba là số nào?
\choice
{\True $45$}
{$44$}
{$46$}
{$43$}
\loigiai{
Thứ tự từ lớn đến bé: $89 > 67 > 45 > 23 > 12$. Số đứng thứ ba là số $45$.
\textbf{Đáp số:} $45$.\par
\textbf{Đáp án của câu hỏi là A}
}
\end{ex}

\begin{ex} %[1D6V22]
Cho các số: $72, 35, 56, 91, 48$. Sắp xếp theo thứ tự từ bé đến lớn, số đứng ở vị trí thứ hai là số nào?
\choice
{$47$}
{$49$}
{$46$}
{\True $48$}
\loigiai{
Thứ tự từ bé đến lớn: $35 < 48 < 56 < 72 < 91$. Số đứng thứ hai là số $48$.
\textbf{Đáp số:} $48$.\par
\textbf{Đáp án của câu hỏi là D}
}
\end{ex}

\begin{ex} %[1D6V22]
Tìm số tròn chục lớn hơn 55 và bé hơn 65.
\choice
{$50$}
{\True $60$}
{$70$}
{$59$}
\loigiai{
Trong khoảng từ 55 đến 65 chỉ có duy nhất số tròn chục $60$.
\textbf{Đáp số:} $60$.\par
\textbf{Đáp án của câu hỏi là B}
}
\end{ex}

\begin{ex} %[1D6V22]
Tìm số tròn chục lớn hơn 72 và bé hơn 88.
\choice
{$70$}
{$90$}
{\True $80$}
{$79$}
\loigiai{
Số tròn chục thỏa mãn là số $80$.
\textbf{Đáp số:} $80$.\par
\textbf{Đáp án của câu hỏi là C}
}
\end{ex}

\begin{ex} %[1D6V22]
Tìm số lớn nhất có hai chữ số mà chữ số hàng chục là 7.
\choice
{\True $79$}
{$78$}
{$80$}
{$77$}
\loigiai{
Để số lớn nhất, chữ số hàng đơn vị phải lớn nhất là 9. Vậy số đó là $79$.
\textbf{Đáp số:} $79$.\par
\textbf{Đáp án của câu hỏi là A}
}
\end{ex}

\begin{ex} %[1D6V22]
Tìm số bé nhất có hai chữ số mà chữ số hàng chục là 5.
\choice
{$40$}
{$60$}
{$49$}
{\True $50$}
\loigiai{
Để số bé nhất, chữ số hàng đơn vị phải bé nhất là 0. Vậy số đó là $50$.
\textbf{Đáp số:} $50$.\par
\textbf{Đáp án của câu hỏi là D}
}
\end{ex}

\begin{ex} %[1D6V22]
Có bao nhiêu số có hai chữ số lớn hơn 85 và bé hơn 90?
\choice
{$3$}
{\True $4$}
{$5$}
{$2$}
\loigiai{
Các số đó là: $86, 87, 88, 89$. Có tất cả $4$ số.
\textbf{Đáp số:} $4$.\par
\textbf{Đáp án của câu hỏi là B}
}
\end{ex}

\begin{ex} %[1D6V22]
Có bao nhiêu số tròn chục lớn hơn 25 và bé hơn 75?
\choice
{$4$}
{$6$}
{\True $5$}
{$3$}
\loigiai{
Các số tròn chục đó là: $30, 40, 50, 60, 70$. Có tất cả $5$ số.
\textbf{Đáp số:} $5$.\par
\textbf{Đáp án của câu hỏi là C}
}
\end{ex}

\begin{ex} %[1D6V22]
Cho ba số: $a = 64, b = 46, c = 66$. Sắp xếp ba số theo thứ tự từ bé đến lớn, số đứng ở đầu tiên là số nào?
\choice
{\True $46$}
{$45$}
{$47$}
{$44$}
\loigiai{
So sánh: $46 < 64 < 66$. Số bé nhất đứng đầu tiên là $46$.
\textbf{Đáp số:} $46$.\par
\textbf{Đáp án của câu hỏi là A}
}
\end{ex}

\textbf{Phần 2. Trắc nghiệm đúng sai}
\setcounterpageref{ex}{0}

\begin{ex} %[1D6V22]
Xét tính đúng hoặc sai của các khẳng định toán học sau đây:
\choiceTF
{Kết quả của bài toán «Lấy số lớn nhất trong các số: $35, 82, 53, 28$ trừ đi số bé nhất trong các số đó được kết quả bằng bao nhiêu» là $55$.}
{Kết quả của bài toán «Tìm số tự nhiên $y$ tròn chục sao cho: $40 < y < 60$» là $51$.}
{\True Kết quả của bài toán «Trong một đợt kiểm tra, bạn An được 85 điểm, bạn Bình được 90 điểm, bạn Cúc được 88 điểm. Bạn nào đạt điểm cao nhất? (An: 1, Bình: 2, Cúc: 3)» là $2$.}
{\True Kết quả của bài toán «Tìm số có hai chữ số giống nhau lớn hơn 70 và bé hơn 80» là $77$.}
\loigiai{
\begin{itemize}
    \item Ý a) Sai: Kết quả đúng là $54$, không phải $55$. Số lớn nhất là 82, số bé nhất là 28: $82 - 28 = 54$.
\textbf{Đáp số:} $54$.
    \item Ý b) Sai: Kết quả đúng là $50$, không phải $51$. Số tròn chục nằm giữa 40 và 60 là số $50$.
\textbf{Đáp số:} $50$.
    \item Ý c) Đúng: So sánh điểm: $90 > 88 > 85$. Bạn Bình đạt điểm cao nhất (ghi số 2).
\textbf{Đáp số:} $2$.
    \item Ý d) Đúng: Số có hai chữ số giống nhau trong khoảng này là số $77$.
\textbf{Đáp số:} $77$.
\end{itemize}
\begin{itemize}
    \item a) \textbf{Sai}.
    \item b) \textbf{Sai}.
    \item c) \textbf{Đúng}.
    \item d) \textbf{Đúng}.
\end{itemize}
}
\end{ex}

\begin{ex} %[1D6V22]
Xét tính đúng hoặc sai của các bài toán và phép tính sau đây:
\choiceTF
{\True Kết quả của bài toán «Tìm số có hai chữ số giống nhau lớn hơn 80 và bé hơn 95» là $88$.}
{Kết quả của bài toán «Hiệu giữa số lớn nhất có hai chữ số và số lớn nhất có một chữ số là bao nhiêu» là $91$.}
{Kết quả của bài toán «Hiệu giữa số bé nhất có hai chữ số và số bé nhất có một chữ số khác 0 là bao nhiêu» là $8$.}
{\True Kết quả của bài toán «Trong một đợt kiểm tra, bạn An được 85 điểm, bạn Bình được 90 điểm, bạn Cúc được 88 điểm. Bạn nào đạt điểm cao nhất? (An: 1, Bình: 2, Cúc: 3)» là $2$.}
\loigiai{
\begin{itemize}
    \item Ý a) Đúng: Số có hai chữ số giống nhau thỏa mãn là số $88$.
\textbf{Đáp số:} $88$.
    \item Ý b) Sai: Kết quả đúng là $90$, không phải $91$. Số lớn nhất có hai chữ số là 99, số lớn nhất có một chữ số là 9: $99 - 9 = 90$.
\textbf{Đáp số:} $90$.
    \item Ý c) Sai: Kết quả đúng là $9$, không phải $8$. Số bé nhất có hai chữ số là 10, số bé nhất có một chữ số khác 0 là 1: $10 - 1 = 9$.
\textbf{Đáp số:} $9$.
    \item Ý d) Đúng: So sánh điểm: $90 > 88 > 85$. Bạn Bình đạt điểm cao nhất (ghi số 2).
\textbf{Đáp số:} $2$.
\end{itemize}
\begin{itemize}
    \item a) \textbf{Đúng}.
    \item b) \textbf{Sai}.
    \item c) \textbf{Sai}.
    \item d) \textbf{Đúng}.
\end{itemize}
}
\end{ex}

\textbf{Phần 3. Trả lời ngắn}
\setcounterpageref{ex}{0}

\begin{ex} %[1D6V22]
Tìm số tự nhiên $x$ có hai chữ số biết: $x$ lớn hơn 97 và bé hơn 99.
\shortans{98}
\loigiai{
\textbf{Đáp số:} $98$\par
Số tự nhiên nằm giữa 97 và 99 là số $98$.
\textbf{Đáp số:} $98$.
}
\end{ex}

\begin{ex} %[1D6V22]
Có bao nhiêu số có hai chữ số có chữ số tận cùng là 0?
\shortans{9}
\loigiai{
\textbf{Đáp số:} $9$\par
Đó là các số tròn chục: $10, 20, 30, \dots, 90$. Có tất cả $9$ số.
\textbf{Đáp số:} $9$.
}
\end{ex}

\begin{ex} %[1D6V22]
Tìm số tròn chục liền trước của số 100.
\shortans{90}
\loigiai{
\textbf{Đáp số:} $90$\par
Các số tròn chục: $80, 90, 100$. Số tròn chục liền trước 100 là $90$.
\textbf{Đáp số:} $90$.
}
\end{ex}

\begin{ex} %[1D6V22]
Hai bạn Mai và Lan sưu tầm tem. Mai có 52 con tem, Lan có 48 con tem. Hỏi Mai nhiều hơn Lan bao nhiêu con tem?
\shortans{4}
\loigiai{
\textbf{Đáp số:} $4$\par
Số tem Mai nhiều hơn Lan là: $52 - 48 = 4$ (con tem).
\textbf{Đáp số:} $4$.
}
\end{ex}

\begin{ex} %[1D6V22]
Bác An thu hoạch được 65 bắp ngô, bác Ba thu hoạch được 70 bắp ngô. Bác Ba thu hoạch nhiều hơn bác An bao nhiêu bắp ngô?
\shortans{5}
\loigiai{
\textbf{Đáp số:} $5$\par
Số bắp ngô bác Ba nhiều hơn bác An là: $70 - 65 = 5$ (bắp ngô).
\textbf{Đáp số:} $5$.
}
\end{ex}

\begin{ex} %[1D6V22]
Cho dãy số: $15, 20, 25, 30, \dots, 40$. Số thích hợp điền vào chỗ chấm là số nào?
\shortans{35}
\loigiai{
\textbf{Đáp số:} $35$\par
Dãy số cách đều 5 đơn vị: $30 + 5 = 35$.
\textbf{Đáp số:} $35$.
}
\end{ex}
